\documentclass[11pt,letterpaper]{article}
\usepackage[utf8]{inputenc}
\usepackage[margin=1in]{geometry}
\usepackage{amsmath}
\usepackage{amsfonts}
\usepackage{booktabs}
\usepackage{longtable}
\usepackage{hyperref}
\usepackage{listings}
\usepackage{xcolor}

\definecolor{codegreen}{rgb}{0,0.6,0}
\definecolor{codegray}{rgb}{0.5,0.5,0.5}
\definecolor{codepurple}{rgb}{0.58,0,0.82}
\definecolor{backcolour}{rgb}{0.95,0.95,0.92}
\definecolor{cautionred}{rgb}{0.55,0,0}

\lstset{
    backgroundcolor=\color{backcolour},
    commentstyle=\color{codegreen},
    keywordstyle=\color{magenta},
    numberstyle=\tiny\color{codegray},
    stringstyle=\color{codepurple},
    basicstyle=\ttfamily\footnotesize,
    breakatwhitespace=false,
    breaklines=true,
    captionpos=b,
    keepspaces=true,
    numbers=left,
    numbersep=5pt,
    showspaces=false,
    showstringspaces=false,
    showtabs=false,
    tabsize=2
}

% Caution box. Deliberately does NOT capture its body as a macro argument,
% because several caution blocks contain lstlisting environments and
% verbatim-like environments cannot be read from inside a captured argument.
\newenvironment{caution}[1]{%
    \par\addvspace{8pt}%
    \noindent{\color{cautionred}\rule{\linewidth}{0.8pt}}\par\vspace{2pt}
    \noindent\textbf{\color{cautionred}Caution: #1}\par\vspace{2pt}
}{%
    \par\vspace{2pt}\noindent{\color{cautionred}\rule{\linewidth}{0.4pt}}%
    \par\addvspace{8pt}%
}

\title{Arctos Robotic Arm --- Motor and CAN Reference Guide}
\author{Consolidated from project debugging logs, source-code inspection, \\ and empirical bench testing}
\date{September 2026}

\begin{document}
\maketitle
\tableofcontents
\newpage

\section{About This Guide}
This is the \emph{reference} document for the Arctos arm's motors and CAN bus: how
the bus works, how the protocol is framed, how to convert encoder counts to joint
angles, and what has actually been verified against hardware.

It is deliberately \textbf{not} a step-by-step procedure. Two companion SOPs cover
that, and each one points back here for background:

\begin{itemize}
    \item \texttt{sop\_nema17.tex} --- bringing up a NEMA 17 joint (MKS Servo42D).
          Joints A, B and C.
    \item \texttt{sop\_nema23.tex} --- bringing up a NEMA 23 joint (MKS\_57D).
          Joint X.
\end{itemize}

This guide supersedes \texttt{arctos\_joint\_b\_guide.tex}, which was removed. All
of its generally-applicable content was folded in here, with several items
\textbf{corrected} in the process --- see Section~\ref{sec:superseded} for exactly
what changed and why, so that nothing is silently re-derived from the old version.

\section{Core Concepts}

\subsection{CAN bus, in one paragraph}
CAN is a differential two-wire multi-drop bus. Every device shares the same pair
(\texttt{CAN\_H} and \texttt{CAN\_L}) plus a common ground, and each device is
addressed by an ID. Every node on the bus must agree on the bitrate; the Arctos arm
runs at \textbf{500~kbit/s}. Because all nodes share one pair, two devices
configured with the same CAN ID will collide and neither will behave correctly.

\subsection{Gearboxes and closed-loop steppers, in one paragraph}
Each joint is a stepper motor driving a reduction gearbox. The reduction multiplies
torque and divides speed and angular error: with a 67.82:1 gearbox, one motor
revolution is only $360/67.82 \approx 5.31^\circ$ at the joint. The MKS controllers
are \emph{closed-loop}: a magnetic encoder on the motor shaft reports true position,
so the controller can detect that it has fallen behind. Two consequences matter
constantly in practice --- the encoder measures the \textbf{motor shaft, not the
joint output}, and a stalled stepper converts its full phase current into heat with
no motion to carry it away.

\section{Hardware}

\subsection{Joint, motor and CAN ID map}
\begin{table}[h]
\centering
\begin{tabular}{lllll}
\toprule
\textbf{Joint} & \textbf{Role} & \textbf{CAN ID} & \textbf{Motor / driver} & \textbf{Gear ratio} \\ \midrule
X & Base/waist rotation & \texttt{0x01} & NEMA 23 / MKS\_57D    & 13.5:1  \\
Y & (not confirmed)     & \texttt{0x02} & ---                   & 150.0:1 \\
Z & (not confirmed)     & \texttt{0x03} & ---                   & 150.0:1 \\
A & Elbow               & \texttt{0x04} & NEMA 17 / Servo42D    & 48.0:1  \\
B & Wrist rotation      & \texttt{0x05} & NEMA 17 / Servo42D    & 67.82:1 \\
C & Wrist joint         & \texttt{0x06} & NEMA 17 / Servo42D    & 67.82:1 \\
\bottomrule
\end{tabular}
\caption{Joint map. CAN IDs and gear ratios match \texttt{arctos\_controller.yaml}.
Joint A's \texttt{0x04} and its 48:1 ratio are empirically verified; the others are
config values and should be confirmed with \texttt{can\_id\_scan.py} before use.}
\end{table}

\begin{caution}{Always confirm a joint's CAN ID against real hardware}
This project's documents have carried wrong CAN IDs for Joint~A and Joint~C at
various times, and Joint~C's motor was found at \texttt{0x05} when both the config
and the docs said otherwise. Run \texttt{can\_id\_scan.py} --- which sends only the
read-only \texttt{0x31} query and never enables coils --- before trusting any table,
including this one.
\end{caution}

\subsection{Motor specifications}
\begin{table}[h]
\centering
\begin{tabular}{lll}
\toprule
\textbf{Parameter} & \textbf{NEMA 17 (Servo42D)} & \textbf{NEMA 23 (MKS\_57D)} \\ \midrule
Holding torque    & 24~N\textperiodcentered cm & 1.8~N\textperiodcentered m \\
Rated current     & 1.2~A                      & 2.8~A \\
Shaft diameter    & 5~mm                       & 6.35~mm \\
Step angle        & 1.8\textdegree\ (200 full steps/rev) & 1.8\textdegree\ (200 full steps/rev) \\
Microstepping     & 16$\times$ (3200 pulses/rev) & 16$\times$ (3200 pulses/rev) \\
Encoder           & 16{,}384 counts/rev (14-bit) & 16{,}384 counts/rev (14-bit) \\
Joints            & A, B, C                    & X \\
\bottomrule
\end{tabular}
\end{table}

\begin{caution}{Do not confuse motor steps, microstep pulses, and encoder counts}
Three different resolutions are in play and mixing them is the single most common
arithmetic error in this project:
\begin{itemize}
    \item \textbf{200} full steps per motor revolution --- the physical stepping resolution.
    \item \textbf{3200} microstep pulses per motor revolution --- the scale that
          \texttt{POSITION\_CONTROL (0xFD)} commands are expressed in.
    \item \textbf{16{,}384} encoder counts per motor revolution --- the scale that
          \texttt{READ\_ENCODER (0x31)} reports in.
\end{itemize}
One microstep pulse is $16384/3200 = 5.12$ encoder counts. Commanded motion and
measured motion are therefore \emph{never} in the same units.
\end{caution}

\subsection{CANable 2 adapter}
\begin{itemize}
    \item USB VID:PID \texttt{16d0:117e}.
    \item Contains its own CAN controller --- \textbf{do not} add an external
          MCP2515 module in series with it.
    \item Only three wires are needed between adapter and controller:
          \texttt{CAN\_H}, \texttt{CAN\_L}, \texttt{GND}. Do \emph{not} connect the
          CANable's power lines to the motor controller.
\end{itemize}

\begin{lstlisting}[caption=Physical signal path]
Ubuntu host
     |
    USB
     |
 CANable 2
     |
CAN_H / CAN_L / GND
     |
MKS controller (Servo42D or 57D)
     |
  Joint motor
\end{lstlisting}

\subsection{Bus termination check (power off)}
With the system powered off, measure resistance between \texttt{CAN\_H} and
\texttt{CAN\_L}:

\begin{table}[h]
\centering
\begin{tabular}{ll}
\toprule
\textbf{Reading} & \textbf{Meaning} \\ \midrule
$\approx 60\,\Omega$  & Correct --- two 120$\Omega$ terminators in parallel \\
$\approx 120\,\Omega$ & Only one terminator present \\
Infinite / open       & Wiring is disconnected \\
Very low              & Possible short \\
\bottomrule
\end{tabular}
\end{table}

\section{The CAN Interface on Ubuntu}
\label{sec:can-interface}

The CANable presents as a serial device (\texttt{slcan}), which is bridged to a
Linux network interface by \texttt{slcand}. This must be rebuilt every time the
adapter is plugged in or re-enumerates.

\begin{lstlisting}[language=bash, caption=Bring-up]
lsusb                                   # expect: ID 16d0:117e CANable2
ls -l /dev/serial/by-id/                # stable name, survives renumbering
sudo pkill slcand
sudo ip link delete can0 2>/dev/null
sudo slcand -o -c -s6 /dev/serial/by-id/<your-canable> can0
sudo ip link set can0 txqueuelen 1000
sudo ip link set up can0
ip -br link show can0                   # expect: can0  UP
\end{lstlisting}

\begin{caution}{The \texttt{slcand} bitrate flag is a coded index, not a number}
\texttt{-s0}=10k, \texttt{-s1}=20k, \texttt{-s2}=50k, \texttt{-s3}=100k,
\texttt{-s4}=125k, \texttt{-s5}=250k, \textbf{\texttt{-s6}=500k}, \texttt{-s7}=800k,
\texttt{-s8}=1M. The Arctos bus is 500~kbit/s, so \texttt{-s6} is correct. Earlier
project scripts used \texttt{-s3}, which silently attached at 100~kbit/s.
\end{caution}

\begin{caution}{On slcan, the bitrate is set at attach time and \texttt{ip} cannot show it}
Do \emph{not} use \texttt{ip link set can0 up type can bitrate 500000}. That netlink
attribute is for native CAN controller drivers and is not supported on slcan
interfaces. Consequently \texttt{ip -details link show can0} reports
\texttt{bitrate 0} and \texttt{clock 0} on a perfectly healthy CANable --- this is
normal and is \textbf{not} a fault to chase. An older version of this guide told the
reader to verify by looking for ``\texttt{bitrate 500000}''; that check can never
pass on this hardware.
\end{caution}

\begin{caution}{Kernel CAN error counters are meaningless on slcan}
The \texttt{re-started / bus-errors / arbit-lost / error-warn / error-pass /
bus-off} counters in \texttt{ip -details -statistics link show can0} are not
populated by the slcan driver --- they read zero regardless of what the bus is
doing. Do not use them as evidence that a bus is healthy, and note that the
``total silence with zero bus errors implies no CAN transceiver'' heuristic
(Section~\ref{sec:rs485}) requires a native CAN controller to be meaningful.
\end{caution}

\begin{lstlisting}[language=bash, caption=Shutdown]
# 1. Stop any running control script (Ctrl+C)
sudo ip link set can0 down
sudo pkill slcand
# then physically unplug the adapter
\end{lstlisting}

\section{Controller Panel Configuration}
\label{sec:panel}

Confirm these on the controller's own display before attempting CAN communication:

\begin{table}[h]
\centering
\begin{tabular}{ll}
\toprule
\textbf{Setting} & \textbf{Value} \\ \midrule
CAN ID                & the joint's ID from the map above \\
CAN Rate              & \texttt{500} (kbit/s) \\
CAN Response (CANRSP) & \texttt{ENABLE} \\
CAN CRC               & \texttt{SUM-8} \\
Working current       & at or below the motor's rated current \\
Working mode          & see caution below \\
\bottomrule
\end{tabular}
\end{table}

Save and power-cycle the controller after any change. The CAN ID in particular must
be committed, not merely displayed.

\begin{caution}{Set the working current to the motor's rating, not the panel default}
A controller was found in this project storing a 3000~mA default against a motor
rated 1.2~A --- roughly 2.5$\times$. Check this on every new controller.
\end{caution}

\begin{caution}{Working-mode mismatch can imitate a communication fault}
The ROS~2 driver defines six modes and defaults to \texttt{CR\_vFOC} (2):
\begin{lstlisting}[language=C++]
enum class MotorMode {
    CR_OPEN = 0, CR_CLOSE = 1, CR_vFOC = 2,
    SR_OPEN = 3, SR_CLOSE = 4, SR_vFOC = 5
};
\end{lstlisting}
If the panel's mode disagrees with what a motion command assumes, the controller can
accept and checksum-acknowledge a frame without rotating the shaft. Confirm the mode
before concluding a ``no motion'' result is a wiring or CAN problem.
\end{caution}

\section{The MKS CAN Protocol}

\subsection{Frame structure and checksum}
Every frame is \texttt{[command byte] [data bytes...] [checksum byte]}, where the
checksum is an 8-bit sum that \textbf{includes the CAN ID}:

\begin{lstlisting}[language=python]
def checksum(motor_id, data_bytes):
    return (motor_id + sum(data_bytes)) & 0xFF
\end{lstlisting}

\begin{table}[h]
\centering
\begin{tabular}{llll}
\toprule
\textbf{Frame (no checksum)} & \textbf{CAN ID} & \textbf{Arithmetic} & \textbf{Checksum} \\ \midrule
\texttt{31}     & \texttt{0x04} & \texttt{0x31+0x04} & \texttt{35} \\
\texttt{31}     & \texttt{0x05} & \texttt{0x31+0x05} & \texttt{36} \\
\texttt{3A}     & \texttt{0x05} & \texttt{0x3A+0x05} & \texttt{3F} \\
\texttt{F3 01}  & \texttt{0x05} & \texttt{0xF3+0x01+0x05} & \texttt{F9} \\
\texttt{F7}     & \texttt{0x06} & \texttt{0xF7+0x06} & \texttt{FD} \\
\bottomrule
\end{tabular}
\caption{Worked checksum examples. Note the checksum differs per joint --- a frame
copied from another joint's documentation will be silently rejected.}
\end{table}

\begin{caution}{Never hand-invent a checksum, and never reuse another joint's frame}
Omitted or incorrect checksum bytes were the direct cause of several ``the motor
ignored my command'' failures in this project. Compute it from the formula every
time, including for the emergency stop.
\end{caution}

\subsection{Command reference}
\begin{longtable}{lll}
\toprule
\textbf{Name} & \textbf{Byte} & \textbf{Status} \\ \midrule
\endhead
\texttt{QUERY\_MOTOR}                & \texttt{0xF1} & Liveness check. Response \texttt{F1 01 <crc>}. \\
\texttt{READ\_ENCODER}               & \texttt{0x31} & \textbf{Verified.} Use this one. See below. \\
\texttt{READ\_ENCODER} (carry form)  & \texttt{0x30} & Superseded --- see Section~\ref{sec:superseded}. \\
\texttt{READ\_VELOCITY}              & \texttt{0x32} & Responds; zero at rest. Layout unconfirmed. \\
\texttt{READ\_ERROR}                 & \texttt{0x39} & Signed 32-bit position deficit. Grows when stalled. \\
\texttt{READ\_ENABLE\_STATE}         & \texttt{0x3A} & \texttt{01}=enabled, \texttt{00}=disabled. Verified. \\
\texttt{READ\_SHAFT\_PROTECTION\_STATE} & \texttt{0x3E} & \texttt{00}=clear. Non-zero latches the driver. \\
\texttt{SET\_WORKING\_MODE}          & \texttt{0x82} & See panel caution. \\
\texttt{ENABLE\_SHAFT\_PROTECTION}   & \texttt{0x88} & \textbf{Avoid} --- false positives, latches. \\
\texttt{ENABLE\_MOTOR}               & \texttt{0xF3} & \texttt{01}=enable, \texttt{00}=disable. Verified. \\
\texttt{RELATIVE\_POSITION}          & \texttt{0xF4} & Unsafe packing --- see motion section. \\
\texttt{ABSOLUTE\_POSITION}          & \texttt{0xF5} & \textbf{Unsafe} --- see motion section. \\
\texttt{SPEED\_CONTROL}              & \texttt{0xF6} & Not exercised. \\
\texttt{EMERGENCY\_STOP}             & \texttt{0xF7} & Verified with checksum. \\
\texttt{POSITION\_CONTROL}           & \texttt{0xFD} & \textbf{Verified. Preferred for all motion.} \\
\bottomrule
\end{longtable}

\subsection{Reading the encoder}
\label{sec:encoder}

Use \texttt{0x31}. The response is
\texttt{[0x31][6 bytes of signed position][checksum]}, where the six bytes are a
big-endian signed 48-bit multi-turn count on the 16{,}384-counts-per-revolution
scale:

\begin{lstlisting}[language=python, caption=Verified position extraction]
raw48 = int.from_bytes(data[1:7], byteorder="big", signed=True)
motor_degrees = raw48 * 360.0 / 16384
joint_degrees = motor_degrees / gear_ratio
\end{lstlisting}

This is multi-turn and does not wrap at one revolution. It is reset to approximately
zero when the controller is power-cycled, so it is a \emph{relative} reference from
power-on, not an absolute machine coordinate.

\begin{caution}{The encoder measures the motor shaft, not the joint output}
Gearbox backlash and any lost motion downstream of the motor are invisible to this
encoder. It cannot be used to measure output-side backlash; that needs a dial
indicator on the output.
\end{caution}

\subsection{Encoder non-linearity}
The magnetic encoder carries a once-per-revolution sinusoidal error caused by the
sensing magnet sitting slightly off-centre on the shaft. Measured on Joint~A over
6{,}725 samples: \textbf{1.465\textdegree{} peak-to-peak at the motor shaft}, about
0.407\% of a revolution.

Two practical consequences:
\begin{itemize}
    \item It cancels over any whole number of motor revolutions. Choosing a step
          size that is an integer number of motor turns makes it vanish --- Joint~A
          steps of exactly 6 motor revolutions each landed at 100.0\% every time,
          where 90\textdegree{} motor-shaft steps scattered 98.7--101.4\%.
    \item Through a reduction it becomes negligible at the joint: 1.465\textdegree{}
          through 48:1 is 0.031\textdegree{}.
\end{itemize}
Repeatability is unaffected --- returning to the same shaft angle repeats to better
than one encoder count. Absolute readings carry the cyclic error; relative and
repeat positioning do not.

\subsection{Motion commands}
\label{sec:motion}

\textbf{Use \texttt{POSITION\_CONTROL (0xFD)}.} It is the command the project's own
working jog code uses, and the only one verified safe here.

\begin{lstlisting}[caption=0xFD frame layout]
[0xFD]
[direction | ((speed >> 8) & 0x0F)]   direction: 0x00 forward, 0x80 reverse
[speed & 0xFF]                        speed is approximately rpm (see below)
[accel]
[(pulses >> 16) & 0xFF]               ALWAYS-POSITIVE magnitude,
[(pulses >>  8) & 0xFF]               on the 3200-pulses-per-motor-rev scale
[pulses & 0xFF]
[checksum]
\end{lstlisting}

\[
\text{pulses} = |\text{joint degrees}| \times \frac{3200 \times \text{gear\_ratio}}{360}
\]

The driver replies with an unsolicited two-frame status sequence:
\texttt{FD 01} (``started'') and then, if the commanded distance was actually
achieved, \texttt{FD 02} (``complete''). \textbf{Absence of \texttt{FD 02} is the
reliable signal that a command did not finish.}

\begin{caution}{Do not use \texttt{0xF5}, and be careful with \texttt{0xF4}}
\texttt{0xF5} packs an absolute target through a plain \texttt{\& 0xFFFFFF}. On a
joint whose raw counter had drifted to roughly $-274{,}000$ counts, a small intended
nudge became an enormous, non-settling, uninterruptible move that required cutting
motor power. \texttt{0xF4} packs a possibly-negative delta the same way. The
\texttt{0xFD} pattern above --- a separate direction byte plus an always-positive
magnitude --- does not have this failure mode.
\end{caution}

\begin{caution}{Never retry a stalled command into an obstruction}
Re-sending the same command to a joint that has not moved does not change the
situation; it only dumps stall current into the motor as heat. Nine consecutive
full-torque retries into a hard stop drove one joint's internal position error to
$-124{,}374$ counts. If a step under-delivers, stop, or back off and re-attempt ---
do not repeat blindly.
\end{caution}

\subsection{Speed and its limits}
The \texttt{speed} field is very close to motor rpm. Measured on a NEMA 17 with no
load:

\begin{table}[h]
\centering
\begin{tabular}{lll}
\toprule
\textbf{Commanded} & \textbf{Actual} & \textbf{Behaviour} \\ \midrule
10--225  & tracks to within 0.5\% & linear; \texttt{rpm} $= 1.001 \times$ \texttt{speed} \\
250      & 98.7\%                 & edge of linear \\
275--325 & 96.5\%--90\%           & rolling off \\
350+     & saturates $\approx$310~rpm & losing steps \\
\bottomrule
\end{tabular}
\end{table}

There is no hard stall --- speed rolls off gracefully and plateaus. \textbf{But
above roughly 250 the motor receives more pulses than it executes, so position is
silently lost.} For anything where position matters, stay at or below 225.

\section{Converting Encoder Counts to Joint Degrees}

\[
\text{counts per joint degree} = \frac{\text{gear\_ratio} \times 16384}{360}
\]

\begin{table}[h]
\centering
\begin{tabular}{lll}
\toprule
\textbf{Gear ratio} & \textbf{Counts per joint degree} & \textbf{Joints} \\ \midrule
13.5:1   & 614.4    & X \\
48.0:1   & 2184.5   & A \\
67.82:1  & 3086.56  & B, C \\
150.0:1  & 6826.7   & Y, Z \\
\bottomrule
\end{tabular}
\end{table}

\begin{caution}{Three wrong conversion constants exist in this project's history}
Do not use any of them. (1) ``$\approx$3300 counts per 90\textdegree'' describes the
\emph{motor shaft}, not the joint --- through 67.82:1 that same motion is only
$\approx 1.07^\circ$ of joint travel. (2) \texttt{COUNTS\_PER\_REV = 16050} --- the
encoder is 14-bit, so 16{,}384 ($2^{14}$) is correct and source-verified.
(3) \texttt{pulses\_per\_degree = -813.0} in \texttt{mks\_driver.py}, off by roughly
3.8$\times$, with an accompanying ``safety cage'' that estimated limits using a
different multiplier than the one actually transmitted --- so it could not reliably
block an unsafe command. Use only the formula above.
\end{caution}

\section{Safety}
\label{sec:safety}

\begin{table}[h]
\centering
\begin{tabular}{lll}
\toprule
\textbf{Action} & \textbf{Frame} & \textbf{Example (Joint A, \texttt{0x04})} \\ \midrule
Enable coils    & \texttt{F3 01 <crc>} & \texttt{cansend can0 004\#F301F8} \\
Disable coils   & \texttt{F3 00 <crc>} & \texttt{cansend can0 004\#F300F7} \\
Emergency stop  & \texttt{F7 <crc>}    & \texttt{cansend can0 004\#F7FB} \\
\bottomrule
\end{tabular}
\caption{Recompute the checksum for the joint you are actually driving.}
\end{table}

\begin{caution}{Compute your joint's e-stop frame before you need it}
A previous version of this documentation gave Joint~C's emergency stop as
\texttt{cansend can0 001\#F7F8} --- addressed to CAN ID \texttt{0x01}. It would have
done nothing to the actual motor. Derive the frame for the joint in front of you,
write it down, and keep it in a second terminal before commanding motion.
\end{caution}

\begin{caution}{Dropping coils on a gravity-loaded joint lets it fall}
Several scripts in this package disable the motor in a \texttt{finally} block. On an
unloaded bench motor that is harmless; on an assembled arm it removes holding
torque. Know which you have before running one.
\end{caution}

\section{Troubleshooting}
\label{sec:troubleshooting}

\begin{longtable}{p{0.34\textwidth}p{0.58\textwidth}}
\toprule
\textbf{Symptom} & \textbf{Likely cause / fix} \\ \midrule
\endhead
\texttt{can0} does not exist &
  \texttt{slcand} died, usually because the adapter re-enumerated. Re-run the
  bring-up in Section~\ref{sec:can-interface}. Check
  \texttt{/dev/serial/by-id/} for the current device. \\
Adapter present but \texttt{can0} missing after replug &
  Expected --- \texttt{slcand} does not restart itself. A udev rule keyed on the
  adapter's serial can automate this. \\
\texttt{ip} shows \texttt{bitrate 0} &
  Normal on slcan. Not a fault. \\
No response from any ID &
  Check the panel's CAN ID, CAN rate, and CANRSP; check \texttt{CAN\_H}/\texttt{CAN\_L}
  and termination; confirm the driver is powered. If still silent, see
  Section~\ref{sec:rs485}. \\
Command acknowledged, motor does not move &
  Distinguish ``not being driven'' from ``driven but cannot move the load'' --- the
  bus cannot tell these apart. Disable coils and turn the joint by hand. \\
Motion for a few degrees, then nothing in either direction &
  If it cannot even \emph{back off}, suspect either a bind stiff enough to absorb
  full torque, or a marginal motor phase connection. Reseating a connector once
  took a joint's travel budget from 1.05\textdegree{} to 3.85\textdegree{}. \\
Gritty resistance by hand, stalls under power &
  Gear mesh needs lubrication. Use a plastic-safe grease (PTFE or silicone).
  Petroleum jelly works as a diagnostic but migrates once warm. \\
Motor warm but not hot &
  Normal. Thermal shutdown makes a motor too hot to touch; ``warm'' is a stepper
  doing its job. \\
Driver stops responding mid-test &
  Check power and connectors first. This has been caused by an accidental unplug
  more than once. \\
\bottomrule
\caption{Failure modes actually observed on this arm.}
\end{longtable}

\subsection{Total silence with a powered driver: the RS485 variant}
\label{sec:rs485}
The MKS SERVO42D/57D ship in both CAN and RS485 variants. The RS485 boards have
\textbf{no CAN transceiver at all} and will never appear on the bus at any bitrate or
address. Joint~X's driver turned out to be the RS485 variant, which is why it has
never communicated. If a joint is silent across a full ID scan and a bitrate sweep,
with power and wiring confirmed, \textbf{check the board's part number before
spending more time in software}.

\section{What This Guide Supersedes}
\label{sec:superseded}

\texttt{arctos\_joint\_b\_guide.tex} was removed and its generally-applicable content
consolidated here. The following items were \textbf{corrected} rather than copied,
so do not re-derive them from any surviving copy of the old document:

\begin{itemize}
    \item \textbf{Encoder opcode.} The old guide standardized on \texttt{0x30} and
          described \texttt{0x31} as unverified. This is backwards: \texttt{0x31} is
          what the real C++ driver uses and what every script in this package now
          uses, with the 48-bit layout given in Section~\ref{sec:encoder}.
    \item \textbf{Bitrate verification.} The old guide instructed the reader to run
          \texttt{ip link set can0 up type can bitrate 500000} and then verify by
          looking for \texttt{bitrate 500000}. Neither works on slcan; the interface
          correctly reports \texttt{bitrate 0}.
    \item \textbf{Bus error counters.} The old troubleshooting advice leaned on
          bus-error counts that the slcan driver never populates.
    \item \textbf{Motion command.} The old guide's safe-first-motion procedure was
          built on \texttt{0xF5}, whose absolute-vs-relative behaviour it correctly
          flagged as unresolved. That question is now moot: use \texttt{0xFD}.
    \item \textbf{Joint-specific framing.} Checksums, e-stop frames and panel
          settings are now presented per-joint rather than hard-coded to Joint~B's
          \texttt{0x05}, which was a copy-paste hazard.
\end{itemize}

\end{document}
